% =============================================================================
% 文件: sections/s1_restate.tex   第一章 问题重述
% =============================================================================
\section{问题重述}

\subsection{问题背景}

干燥是中药材产地加工与饮片生产中的关键工序。热风烘干因设备简单、适应性强而被广泛
采用，其工艺过程通常划分为\textbf{预热平衡阶段}与\textbf{恒温干燥阶段}：前者物料由室温
逐步升温、表面水分开始蒸发，后者在稳定的温湿环境下依靠内部水分向外迁移持续失水。
工艺参数（风温、湿度、风速）选取不当会造成干燥效率低、能耗高、有效成分损失大等后果，
而单纯依靠试验优化存在成本高、周期长的困难，因此需要建立可计算的传热传质模型来
刻画干燥规律并预测干燥时长。

本题将药材简化为长圆柱体：长 $L=25\ \mathrm{cm}$、初始半径 $\Rzero=2\ \mathrm{cm}$，
初始温度 $T_0=28\ ^\circ\mathrm{C}$，初始干基含水率 $C_0=2.55\ \mathrm{kg/kg}$。
烘房的温度与水分浓度随时间变化的数据见附件 1；药材在干燥过程中的半径变化见附件 2；
预热平衡阶段的物性参数见附录 2，恒温干燥阶段与考虑收缩时的经验公式分别见附录 3、附录 4。

\subsection{四个问题的可量化子任务}

把题面逐条转写为“输入—处理—输出—判定标准”的形式，得到如下子任务清单。

\paragraph{问题 1（预热平衡阶段，$0\sim1800\ \mathrm{s}$）}
\begin{itemize}[leftmargin=2.2em,itemsep=0.15em]
  \item \textbf{输入}：附件 1 的烘房温度与水分浓度（$0\sim14400\ \mathrm{s}$，间隔 $60\ \mathrm{s}$，
        共 241 个点）；附录 2 的常数物性（$\rho=820$、$c_p=2600$、$k=0.36$、
        $h=25$、$h_m=8\times10^{-7}$）与含水率依赖的扩散系数
        $D=7\times10^{-9}\mathrm{e}^{-0.89/C}$；初始条件 $T_0=28\ ^\circ\mathrm{C}$、$C_0=2.55$。
  \item \textbf{输出}：（i）表 1、表 2 —— 在 $t=100,300,600,900,1200,1500,1800\ \mathrm{s}$
        与到中心距离 $r=0,0.5,1,1.5,2\ \mathrm{cm}$ 的网格上给出温度与含水率；
        （ii）\texttt{result1.xlsx} —— $0\sim1800\ \mathrm{s}$ 内每 $1\ \mathrm{s}$、
        每 $0.1\ \mathrm{cm}$ 的完整场。
  \item \textbf{判定标准}：结果保留四位小数；$r=2\ \mathrm{cm}$ 列必须是药材\textbf{表面}
        的值而不是最外层计算单元中心的值；温度应单调升温且不超过环境温度，
        含水率应单调下降且非负。
\end{itemize}

\paragraph{问题 2（整个烘干过程的建模，并给出 $3\ \mathrm{h}$ 内结果）}
\begin{itemize}[leftmargin=2.2em,itemsep=0.15em]
  \item \textbf{输入}：同问题 1，但全部经验公式改用附录 3（$\rho$、$c_p$、$k$ 均为含水率的
        函数，$D$ 同时依赖含水率与\textbf{热力学温度}）。
  \item \textbf{输出}：表 3、表 4 —— $0.5,1.0,\dots,3.0\ \mathrm{h}$ 与
        $r=0,0.5,1,1.5,2\ \mathrm{cm}$ 上的温度与含水率；\texttt{result2.xlsx} 为
        每 $1\ \mathrm{s}$、每 $0.1\ \mathrm{cm}$ 的完整场。
  \item \textbf{判定标准}：模型必须覆盖“整个烘干过程”而不是只有 $3\ \mathrm{h}$，
        因此必须显式给出 $t>14400\ \mathrm{s}$（附件 1 数据终点）以后的环境边界假设。
\end{itemize}

\paragraph{问题 3（烘干时长）}
\begin{itemize}[leftmargin=2.2em,itemsep=0.15em]
  \item \textbf{输入}：同问题 2 的模型与参数。
  \item \textbf{输出}：烘干所需时间 $t_3$（单位 h）；表 5 —— 每 $6\ \mathrm{h}$、
        每 $0.5\ \mathrm{cm}$ 的含水率，末行为“烘干结束时间”；
        \texttt{result3.xlsx} 为每 $60\ \mathrm{s}$、每 $0.1\ \mathrm{cm}$ 的完整场。
  \item \textbf{判定标准}：题面要求“各处含水率低于 $0.15\ \mathrm{kg/kg}$”，
        因此终点判据必须是\textbf{局部最大含水率} $\max_r C(r,t)<0.15$，
        而不是截面平均值 —— 两者可相差数小时。
\end{itemize}

\paragraph{问题 4（考虑尺寸变化的烘干时长）}
\begin{itemize}[leftmargin=2.2em,itemsep=0.15em]
  \item \textbf{输入}：附录 4 的经验公式（$\rho=760+90C$、$c_p=1850+2150C/(C+1)$、
        $k=0.12+0.20C/(C+1)$、$D=4.2\times10^{-4}\mathrm{e}^{-0.30/C}\mathrm{e}^{-3850/T}$）；
        附件 2 的半径历史 $R(t)$（$0\sim259200\ \mathrm{s}$，间隔 $1800\ \mathrm{s}$，
        由 $2.000\ \mathrm{cm}$ 单调收缩至 $1.198\ \mathrm{cm}$，体积收缩 $64.1\%$）。
  \item \textbf{输出}：烘干时长 $t_4$；表 6 —— 每 $6\ \mathrm{h}$、每 $0.5\ \mathrm{cm}$
        至药材表面的含水率；\texttt{result4.xlsx} 为每 $60\ \mathrm{s}$、
        每 $0.1\ \mathrm{cm}$ 的完整场（最后一列为“药材表面”）。
  \item \textbf{判定标准}：收缩使半径持续变小，输出列的定义必须与时间相关的 $R(t)$
        自洽；“药材表面”列应取 $r=R(t)$ 处的值。
\end{itemize}

四个问题的共同判定要求是：结果保留四位小数、给出可复现的计算配置
（网格数 $N$、时间步 $\Delta t$、边界处理方式、终点判据），并且必须量化离散误差 ——
脱离离散配置谈“误差很小”是没有意义的。
